%! TEX root = riset_senja_lama.tex
% the main TeX file which is intended to compile, :VimtexReload after adjustment
% :h vimtex-tex-root
\documentclass[../riset_senja.tex]{subfiles}

\begin{document}

\section{Exploring the Place Identity of Public Waterfront Spaces in Parepare City}%
\label{sec:Exploring the Place Identity of Public Waterfront Spaces in Parepare City}

\subsection{Research Background}%
\label{sub:ResearchBackground}

Public space has gained a lot of attention in city planning in recent years, including Public Waterfront Space (PWS). It may be due to the increasing of population that demand a place for recreation.
This population increase is characterized by large-scale housing developments, high traffic density, and the expansion of the business and economic sectors. Therefore people want to escape those by go to the public spaces.

Public space is a place that accommodates all individual or group activities in seeking pleasure and other positive things \citep{sadat2012}.
\cite{sadat2012} explained that public spaces include parks, town squares, waterfronts, and streets. Furthermore, according to its functions, public spaces can be a place for business, recreation, socialization, and cultural activity.

The function of this space has led people to consider it when seeking an accessible place. They use this space for daily activities. For example, firstly, the streets that they traverse. Secondly, the town square for family or friends' recreation. These uses underscore the significance of public space in people's lives.

According to \cite{hartanti2014}, good public space has a strong identity. It is the basis for recognizing something as a separate entity and shaped by characters that exist within it \citep{hartanti2014}.
Additionally, from the dictionary of Webster's Ninth New Collegiate, identity is defined as the distinguishing character or state of a person or thing.
Those characters encompass physical element, activity, social and even meaning \citep{hartanti2014, kaymaz2013}.

A strong identity on public space will benefit significantly to the successful of public space.
For example, with strong identity, people will be interested in coming and staying for longer period \citep{oktay2002}. One of the city that has created strong identity is a Singapore, with the characteristic of lion head in its popular Merlion public space.

Parepare has several public spaces as well, mostly located in near the waterfront. Few of them have just undergone significant development. As a result, it has changed the social, cultural, and morphological order of a spatial composition at the regional and city levels \citep{oktay2002, kaymaz2013}.
According to \citep{Kaymaz13}, the changes or development that have been made will raise a question: has the identity of the space changed or eroded?

The research on public space identity has been conducted numerous times \citep{oktay2015}. However, only a few studies have been carried out in waterfront to make this place more successful.
This public space is more valuable because it has distinguished characteristics gained from its proximity to the water. Also, its significant contribution towards the city development.

\cite{hussein2014} argued that the improved Public Waterfront Space (PWS) will attract numerous visitor to the coastal area.
Waterfront is an interesting and flexible place \citep{hussein2014}. It becomes a sign of a successful city when developed successfully.
Therefore, many cities located near the sea develop PWS, include Parepare. \cite{oktay2002} asserts that the development is a dynamic urban process.

Parepare, nowadays, focuses on developing tourism city \citep{faniapriani2018}.
For example, developing PWS such as Tonrangeng Riverside, Floating Mosque, Mattirotasi Park, Parepare Beach, Cempae Park, and others. The five PWS mentioned are the subject of this research.

Some public spaces have a strong identity in several cities.
For example, first, the old town in Northern Cyprus is characterized by integrated districts, streets with open spaces, and natural landscapes including palm trees and orange gardens \citep{oktay2002}.
Second, some areas in the city of Bogor that are characterized by old trees, colonial buildings, and modern markets \citep{hartanti2014}. Third, Merlion park in Singapore characterized by lion head as mention previously.

The example of those public spaces show that every public spaces own distinguish characters \citep{ahmadi2009}.
Therefore, in order to understand identity of these five waterfront, this research will explore their identity.
This study will answer several research questions as follows:

\begin{enumerate}
	\item What identity is formed in the PWS of Parepare?
	\item How does the people's perception towards the identity of PWS in represent the city?
\end{enumerate}

\subsection{Objective}%
\label{sub:Objective}

The aim of this study is to understand the identity of a number of PWS. Therefore understanding their characteristics is important, which can be explain by the context of spatial elements. The researcher hope that the finding of this research can be used by decision-maker or government to plan the city of Parepare. Those will create a dynamic process of a city that present a successful PWS that have a strong identity, in order to attract many visitors.

\begin{enumerate}
	\item Analyzing the identity of the public coastal space in the city of Parepare.
	\item Understanding the community's perception of the identity of public space in representing the city.
\end{enumerate}


\subsection{Method}%
\label{sub:Method}

\subsubsection{Approach}%
\label{ssub:Approach}

This study aims to analyze the identity of the public spaces along the coast of the city of Parepare. Hence, this research will utilize the research approach to reach its objective.
The research approach is the process of research that includes detailed steps, methods in data collection, analysis, and interpretation of findings \citep{creswell2016}. This research approach is qualitative and quantitative (mixed-method).

The use of a mixed-method approach is better than just one approach in order to address the problem \citep{hartanti2019}. It will have its own part in this research. For qualitative approach, it will produces a comprehensive picture of a phenomenon through direct understanding and observation of the phenomenon in the research object accompanied by data sources \citep{creswell2016}.
The phenomenon consists of physical elements, activities, functions, and others that exist in public spaces in shaping their identity.


According to \cite{groat2013}, qualitative approach emphasizes on the inductive process, with \citeauthor{creswell2016} illustrating this through the use of open-ended questions. The question will be about the characters that visitors consider significant in shaping the identity of the coastal public space.
For example, \cite{hartanti2019} found that old colonial buildings became the identity in one of the public spaces in the city of Bogor.
\cite{creswell2016} explained that this question can evolve as the researcher's understanding of the problem in the field grows \citep{creswell2016}.

While quantitative approach is used to study a specific population or sample.
Data collection is based on research variables. Then, a questionnaire survey is given to find visible characteristics for users in PWS.
Data analysis uses an analytical descriptive design with a cross-sectional approach aimed at assessing how often (frequency) interesting variables (characteristics) occur in a specific space. The analysis uses the R programming language.
Therefore, this research can understand the identity of PWS through their formative characteristics for better design.

\subsubsection{Location}%
\label{ssub:Location}

The location of this research object are 5 (five) selected waterfront public space in Parepare (see \ref{fig:objprop}).
Parepare is located in South Sulawesi Province, about 150 km from the province capital city of Makassar.
The PWS spread across the Parepare seaside, they are listed as follow:

The locations of the research objects are five selected  public waterfront spaces in Parepare (see \ref{fig:objprop}).
Parepare is located in the South Sulawesi Province, approximately 150 km from the provincial capital city of Makassar.
The waterfront public spaces selected spread across the Parepare seaside, and they are listed as follows:

\begin{itemize}
	\item \textit{Tonrangeng Riverside} (The southernmost area in the city of Parepare)
	\item  Terapung Mosque
	\item Mattirotasi Park
	\item Parepare Beach (space in the city center)
	\item  Cempae bridge (the northernmost space in the city)
\end{itemize}

\begin{figure}[!h]
	\begin{center}
		\includegraphics[width=6cm]{objprop}
	\end{center}
	\caption{Parepare city location}
	\label{fig:objprop}
\end{figure}

Initially, this research will conduct observations of the research object in order to collect data for evaluation such as photos, physical measurements, public space layout, and existing characters.
The object evaluation focuses on gathering areas, entrances and exits, parking areas, sidewalks, and public transportation locations in the research object.

\begin{figure}[htb]
	\centering
	\begin{subfigure}{7cm}
		\centering\includegraphics[width=6cm]{tonrangeng}
		\caption{Tonrangeng Riverside}
	\end{subfigure}%
	\begin{subfigure}{7cm}
		\centering\includegraphics[width=6cm]{masjidterap}
		\caption{Terapung Mosque}
	\end{subfigure}\vspace{10pt}
	\begin{subfigure}{7cm}
		\centering\includegraphics[width=6cm]{tamanmatras}
		\caption{Mattirotasi Park}
	\end{subfigure}%
	\begin{subfigure}{7cm}
		\centering\includegraphics[width=6cm]{anjungancem}
		\caption{Cempae Bridge}
	\end{subfigure}\vspace{10pt}
	\begin{subfigure}{7cm}
		\centering\includegraphics[width=6cm]{pantaipare}
		\caption{Parepare Beach}
	\end{subfigure}%

	\label{fig:obj1}
\end{figure}

The research will be conducted on weekdays and holidays. The selected time range is from 6 to 10 in the morning and 6 to 9 in the evening for 9 months in 2025. The reason for choosing these research times is that this time shows the busyness at the research site. This research schedule is included in the table \ref{tab:linimasa}.

The research will take place on everyday.
The selected time range is from 6 to 10 in the morning and 6 to 9 in the evening for a duration of 9 months in 2025. The reason for choosing these research times is that they reflect the busyness at the research site. This research schedule is included in Table \ref{tab:linimasa}.

\begin{landscape}
	\setlength\tabcolsep{1pt}
	\begin{table}
		\centering
		\caption{Research Timeline Plan}
		\label{tab:linimasa}
		\begin{tabular}{|c|p{120pt}|*{\NumMonths}{c|}} \hline
			\multirow{2}{*}{\#} & \multirow{2}{*}{Tasks} \YearHeadRow \\
			\YearMonthSep{2}
			                    & \MonthHeadRow                       \\ \hline
			%%%%%%%%%%%%
			\NewGroup
			\NewBlock{blue}{1}{\NumMonths}
			\Task{Literature Review}
			%%%%%%   1  %%%%%%%
			\NewGroup
			\NewBlock{yellow}{1}{1}
			\Task{Research Methodology}
			\NewBlock{yellow}{2}{2}
			\Task{Public Waterfront Space Characters}
			\NewGroup
			\NewBlock{red}{3}{4}
			\Task{Examination}
			\NewGroup
			\NewBlock{green}{5}{7}
			\Task{Research Proposal Revision}
			%%%%%%   6   %%%%%%% 4
			\NewGroup
			\NewBlock{gray}{5}{10}
			\Task{Preliminary Study}
			\NewBlock{gray}{11}{13}
			\Task{Follow up Data}
			%%%%%%   10    %%%%%%% 5
			\NewGroup
			\NewBlock{orange}{11}{17}
			\Task{Statistic Analysis}
			\NewGroup
			\NewBlock{blue}{18}{25}
			\Task{Writing discussion draft}
			%      \Row{Publication 01}
			%      \Row{Publication 02}
			%      \Row{Publication 03}
			%      \Row{Publication 04}
			\NewGroup
			\NewBlock{red}{26}{28}
			\Task{Examination}
		\end{tabular}
	\end{table}
\end{landscape}

The primary data is the perceptions of visitors towards the characteristics of each public waterfront space obtained through survey questionnaires and observations.
The data source includes respondents in each public waterfront space, specifically those who are willing to respond to each questionnaire and participate in interviews.

The questionnaire questions are arranged based on the objectives of this research. Therefore, the questions refer to the research variables, namely physical characteristics, activities, functions, meanings, and others.
In the beginning, a number of questions explore perception data related to each PWS in terms of comfort, visual quality, distinctiveness, legibility, and sense of place.
Then, some questions explore the characteristics of the respondents. After that, interview questions explore knowledge related to PWS comprehensively.
For example the history, symbolism, and functionality of waterfront areas.

This study uses probability sampling technique with stratified random sampling method. The population used is the total population of Parepare city in 2019, which is 145,178 people \citep{bps2020}.
Therefore, this study uses the Slovin formula to determine the sample size from population by setting the precision value at 89\% as follows:

$$ n= \frac{N}{N(d)^2+1} $$

\begin{tabular}{ll}
	\small
	Description: & n = Number of samples     \\
	             & N = Number of  population \\
	             & d = Precision value 89\%  \\
\end{tabular}

Based on the formula above, the total number of samples is:

$$ n= \frac{145178}{145178(0.11)^2+1} $$
$$ = 82.6 \approx  85\, sample $$


\section{Benefit}%
\label{sec:Benefit}

The city rapid development can change or erode the identity of places \citep{hartanti2012}. This will cause
If this happens, people may no longer recognize a place as something different. Therefore, strong identity is important to attract and support community.
This study will determine the characteristics of each waterfront public space mentioned earlier. The researcher expects that these findings will be used by decision-makers to make effective policies for the future planning of the city.

\bibliographystyle{apalike}
\bibliography{../propid.bib}

\end{document}
